% ==================== 第九节 投资者保护 ======================================
% 股利分配政策 → 滚存利润安排 → 摊薄即期回报与填补措施 → 稳定股价预案
% → 重要承诺汇总（锁定/减持/回购/赔偿）→ 投资者沟通渠道。
\gssection{投资者保护}

\subsection{本次发行后的股利分配政策}

根据《公司章程（草案）》，本次发行上市后，公司实行持续、稳定的利润分配政策，
重视对投资者的合理回报并兼顾公司的可持续发展：

\subsubsection{利润分配的形式及现金分红条件}

公司可以采取现金、股票或者现金与股票相结合的方式分配利润，具备现金分红条件
的，应当优先采用现金分红。公司当年盈利、累计未分配利润为正且现金流满足正常
经营和长期发展需要的，应当进行现金分红。

\subsubsection{现金分红的最低比例}

公司每年以现金方式分配的利润不少于当年实现的可供分配利润的10\%，最近三年以
现金方式累计分配的利润不少于最近三年实现的年均可分配利润的30\%。公司董事会
应当综合考虑所处行业特点、发展阶段、自身经营模式、盈利水平以及是否有重大资
金支出安排等因素，提出差异化的现金分红政策。

\subsubsection{利润分配的决策程序}

公司利润分配方案由董事会结合盈利情况、资金需求拟定，经独立董事发表明确意见
后提交股东会审议批准；调整利润分配政策的，应当经出席股东会的股东所持表决
权的三分之二以上通过。

\subsection{本次发行前滚存利润的分配安排}

经公司2026年第一次临时股东会审议通过，本次发行完成前公司滚存的未分配利润，
由本次发行完成后的新老股东按发行后的持股比例共同享有。

\subsection{本次发行摊薄即期回报及填补措施}

本次发行完成后，公司总股本及净资产规模将有所增加，由于募集资金投资项目产生
效益需要一定周期，公司存在发行当年每股收益、加权平均净资产收益率较上年下降
的风险。为此，公司将通过加快募投项目实施进度、提升经营管理效率、完善利润分
配政策等措施填补被摊薄的即期回报。公司董事、高级管理人员已就切实履行填补回
报措施作出承诺，承诺不无偿或以不公平条件向其他单位或个人输送利益，并将个人
薪酬制度、拟公布的股权激励行权条件与公司填补回报措施的执行情况相挂钩。

\subsection{上市后三年内稳定股价的预案}

公司上市后三年内，若公司股票连续20个交易日的收盘价均低于公司最近一期经审计
的每股净资产（因除权除息等事项调整后），且非因不可抗力因素所致，公司及相关
主体将按照以下顺序启动稳定股价措施：（一）公司回购股份；（二）控股股东、实
际控制人增持公司股份；（三）董事（不含独立董事）、高级管理人员增持公司股份。
具体启动条件、实施程序及单次实施规模以公司股东会审议通过的稳定股价预案为准。

\subsection{相关责任主体作出的重要承诺}

公司及相关责任主体就投资者保护事项作出的各项承诺的完整内容，详见本招股说明
书``附件二、与投资者保护相关的承诺''；主要承诺概要如下：

\subsubsection{股份限售承诺}

控股股东示例控股、实际控制人林××承诺：自公司股票上市之日起36个月内，不转
让或者委托他人管理其直接和间接持有的公司本次发行前已发行的股份，也不由公司
回购该部分股份；公司上市后6个月内如公司股票连续20个交易日的收盘价均低于发
行价，或者上市后6个月期末收盘价低于发行价，其所持股份的锁定期限自动延长6个
月。

公司其他股东按照《公司法》及上交所相关规则履行股份锁定义务；担任董事、高级
管理人员的股东同时承诺，在任职期间每年转让的股份不超过其所持公司股份总
数的25\%，离职后半年内不转让其所持公司股份。

\subsubsection{持股意向及减持承诺}

控股股东、实际控制人承诺：锁定期届满后减持股份的，将遵守中国证监会及上交所
关于股份减持的相关规定，提前披露减持计划，减持价格不低于本次发行的发行价
（因除权除息等事项调整后）。

\subsubsection{欺诈发行的股份购回承诺}

公司及控股股东、实际控制人承诺：若公司不符合发行上市条件，以欺骗手段骗取发
行注册并已经发行上市的，将依法购回本次公开发行的全部新股。

\subsubsection{保荐人先行赔付承诺}

保荐人示例证券承诺：因其为发行人本次公开发行制作、出具的文件有虚假记载、误
导性陈述或者重大遗漏，给投资者造成损失的，将先行赔偿投资者损失。

\subsubsection{未能履行承诺的约束措施}

公司及相关责任主体已就未能履行上述承诺时的约束措施作出安排，包括及时披露未
能履行的具体原因并向投资者道歉、暂不领取（或降低）薪酬或津贴、暂停转让所持
公司股份、依法赔偿投资者损失等。

\subsection{投资者与发行人的沟通渠道}

公司指定董事会办公室为投资者关系管理部门，董事会秘书为投资者关系管理负责人。
投资者可通过公司电话、电子信箱（ir@example.com）及上交所投资者互动平台与公
司沟通。公司信息披露文件刊载于上交所网站（www.sse.com.cn），投资者应以公司
在指定信息披露媒体刊登的公告为准。
